% File: Tables/table_multicase_electrical_summary.tex
% Usage: \MultiCaseElectricalSummaryTable[0.94\textwidth]
% Energies are obtained by trapezoidal integration of the exported power
% profiles over the complete simulated interval.

\newcommand{\MultiCaseElectricalSummaryTable}[1][0.94\textwidth]{%
\begin{table}[H]
    \centering
    \caption{Integrated solar generation and selected load demand for the
    nominal, hot and cold transient cases.}
    \label{tab:multicase_electrical_summary}
    \footnotesize
    \renewcommand{\arraystretch}{1.18}
    \setlength{\tabcolsep}{4.2pt}
    \resizebox{#1}{!}{%
    \begin{tabular}{|l|c|c|c|c|c|}
        \hline
        \textbf{Case} &
        \textbf{Generated energy [Wh]} &
        \textbf{Load energy [Wh]} &
        \textbf{Direct surplus [Wh]} &
        \textbf{Direct deficit [Wh]} &
        \textbf{Net balance [Wh]} \\ \hline
        Nominal & 328.07 & 25.58 & 313.78 & 11.29 & 302.49 \\ \hline
        Hot     & 421.60 & 33.75 & 394.46 & 6.61  & 387.85 \\ \hline
        Cold    & 306.41 & 6.92  & 302.73 & 3.24  & 299.49 \\ \hline
    \end{tabular}%
    }
    \par\vspace{0.3em}
    \begin{minipage}{#1}
        \footnotesize
        \textit{Note:} Direct surplus and deficit are the time integrals of
        the positive and negative parts, respectively, of
        \(Q_{S,\mathrm{elec}}-Q_{\mathrm{load}}\), before applying the battery
        operating limits.
    \end{minipage}
\end{table}%
}
